\section{结论}

本文做出了两项贡献。第一项是解释了为何大模型的强化学习训练对训练-推理不匹配（TIM）如此敏感。在TIM条件下，策略梯度更新包含一个漂移项，该漂移项起到了向采样器进行蒸馏的作用。由于采样器是训练器的一个有偏副本，并且周期性地同步回训练器，这种偏差会在反馈环中不断累积：TIM越严重，训练越早崩溃。先前的工作已经观察到这种偏差会随着训练过程不断累积 \citep{probing, mai_1}。我们证明，仅通过消除漂移项（无需重要性比例），即可在严重量化情况下稳定训练，这确立了漂移项是导致不稳定的根源。这也解释了为何从一个完全不同的模型进行离线蒸馏是稳定的，而在线强化学习却在微小数值差异下崩溃：关键并不在于不匹配的大小，而在于教师模型是否固定，还是追踪学生模型。

我们的第二个贡献是得分中心化，这是一种实用的方法，通过减去期望得分来消除漂移。它是一个无需超参数的加性修正项，可表示为标量损失 \parenCref{eq:sc_loss_main}，且能与重要性采样方法（如 TIS 和 MIS）组合使用。确切的得分中心化需要采样器完整的下一个 token 分布；因此，我们引入了一种近似方法，仅需采样器的 top-$k$  logprob 即可，且在所有实验中表现均与确切得分中心化相当 (\Cref{sec:topk_ablation})。在温和的 TIM 条件下，得分中心化与重要性采样方法表现相当；在严重量化情况下，它是唯一能够稳定训练的方法；而在严重延迟情况下，其与 TIS 或 MIS 的组合表现最佳。

\paragraph{局限性。} 得分中心化可以消除漂移，但剩余的更新衡量的是采样器而非训练器下的奖励与得分之间的协方差 \parenCref{eq:sc_bias}。在严重延迟的情况下，这种不匹配尤为关键，此时得分中心化与重要性采样结合时表现最佳 (\Cref{sec:results_quant_stale})。此外，我们的核心结果来自对短序列故意设置的严重不匹配，作为在较轻不匹配下进行长时间训练的代理。